\documentclass[11pt]{article}
\usepackage[margin=1in]{geometry}
\usepackage[T1]{fontenc}
\usepackage{lmodern}
\usepackage{amsmath}
\usepackage[colorlinks=true,linkcolor=blue,citecolor=blue,urlcolor=blue]{hyperref}
\usepackage[acronym,toc,nopostdot]{glossaries}
\usepackage[nameinlink,noabbrev]{cleveref}

\makenoidxglossaries
\let\printindexedglossaries\printglossaries
\let\printglossaries\printnoidxglossaries

\newglossaryentry{calibration}{
  name={calibration},
  description={the agreement between predicted probabilities and observed frequencies within comparable groups}
}

\newglossaryentry{threshold}{
  name={threshold},
  description={a decision boundary applied to a continuous score}
}

\newglossaryentry{loss}{
  name={loss},
  description={a scalar objective that penalizes prediction errors and unstable classifications}
}

\newacronym{ate}{ATE}{average treatment effect}
\newacronym{roc}{ROC}{receiver operating characteristic}
\newacronym{dgp}{DGP}{data-generating process}

\title{Glossary and Acronym Conversion Stress Paper}
\author{Stage Two Corpus}
\date{}

\begin{document}
\maketitle

\section{Terminology}

The first use of an acronym should expand in running text, as in \gls{ate}, while a later use such as \gls{ate} should use the shortened form. The same pattern appears for \gls{roc} curves and later \gls{roc} diagnostics. This behavior is sensitive to package state because the first-use expansion is computed during compilation.

Glossary terms have a related conversion challenge. A model may use a \gls{threshold} for classification, then revise the \gls{threshold} after observing drift in the \gls{dgp}. In contrast, a \gls{loss} function remains a mathematical object, for example
\[
  \ell(\hat{p},y)= -y\log(\hat{p})-(1-y)\log(1-\hat{p}).
\]
The glossary entry for \gls{calibration} should be listed with its definition, and subsequent mentions of \gls{calibration} should remain concise in the prose.

\section{Reference Use}

The paper deliberately invokes each entry more than once. It uses \gls{dgp} in the first sentence of this section and then refers to the \gls{dgp} again after defining the relevant score. The generated glossary and acronym lists below are part of the expected output surface.

\printglossaries

\end{document}
